% =============================================================================
%  text/HISTORICAL BACKGROUND/ruler-power-dynamics.tex
%  Documents temporal variation in ruler power security; references a figure.
%  统治者权力动态示例：引用图片。
% =============================================================================

We organize the case narrative around three periods that differ in the
security of the ruler's authority. In the first period, the ruler commanded
unrivalled prestige following the founding of the regime, and authority over
both party and state was consolidated. In the second period, policy failures
eroded the ruler's standing, and senior elites mounted increasingly open
challenges to his leadership. In the third period, the ruler reasserted
authority through large-scale reshuffles of the elite.

% 中文对照（可选）
% 本文围绕统治者权力稳固程度不同的三个阶段组织案例叙述：建国初期权力
% 巩固；政策失败后权威受损、精英公开挑战；随后统治者通过大规模人事调整
% 重新确立权威。

Figure~\ref{fig:timeline} provides a stylized illustration of the
power-security dynamics across these periods. The figure is generated from a
standalone TikZ source file (\texttt{pic-tex/example-tikz/example-tikz.tex})
and included as a compiled PDF, following the project's figure workflow.

% 图 1：示意时间线（TikZ 独立编译后插入）
\begin{figure}[htbp]
  \centering
  \includegraphics[width=0.85\textwidth]{example-figure.pdf}
  \caption{Stylized Timeline of Ruler Power Security Across Three Periods}
  \label{fig:timeline}
\end{figure}

This temporal variation provides the identification strategy exploited in the
research design (Section~\ref{sec:sample}): the same ruler, confronting the
same elite network, adopted different strategies under different conditions of
power security.
